% ==============================================================================
% T0 Theory / FFGFT: Document 286 (chapter content, EN)
% Title: The dynamic sector -- kinetic energy as a xi-fractal fraction of the
%        static (T~ * m = 1)
% Author: Johann Pascher
% Date: 21 June 2026
% References: Doc. 047, 067, 080, 204, 255, 263, 285.
% Scripts: kinetik_xi_anteil.py
% Character: problem-statement / research-programme document. It records what is
%   fixed, the single data point, the narrowing, and the structural lead -- and
%   marks the open gap explicitly. No claimed derivation.
% ==============================================================================

\section*{Abstract}
Through the duality $\tilde T\cdot m=1$, FFGFT separates two energy types: \emph{static} (rest mass
$=$ topologically frozen winding, the crystallographic Z$_3$ core) and \emph{kinetic} (free; for
massless states the only form, $E=hf$). The static sector is derived and computed (Z$_3$ mass
operator $+$ mass circle; Doc.~285, BD12-XCORE). The dynamic sector is \emph{not}: the derived rule
that computes the kinetic part, via the duality, as a $\xi$-fractal fraction of the static one is
missing. This document is a \emph{problem-statement document} -- it records what is fixed (the
duality places the window in the kinetic sector), the single existing data point (neutrino,
$m_\nu\sim\xi^2/2\,m_e$, as an ansatz), the narrowing of the $\xi$-power to $p=2$, and the structural
lead $m_{\text{window}}=\xi\cdot\omega$. What remains to be derived is marked explicitly as open. Prefixed is the fixed, reversible Hilbert-space foundation (Doc.~282/270/230) the sector sits on; its actual aim is to find, out of the Hilbert space, the complementary description that maps via $D_f=3-\xi$ and $\xi$ onto SI units.

% ==============================================================================
\section{Where the question comes from}
% ==============================================================================
In Doc.~255, $\tilde T\cdot m=1$ enforces two energy types. The rest mass is \emph{static}: a
topologically stable winding number on $T^4$, fixed in the compact picture by the Z$_3$ mass operator
($r=\sqrt2$, $\theta=2/9$). The kinetic energy is \emph{free}; for a massless particle it is the only
form ($E=hf$). \enquote{Mass is frozen kinetic energy.}

Doc.~285 shows that the empty recursion windows ($3'$, $3''$, \dots) carry no \emph{new} masses but
structure; and that a \emph{dynamics} filling a window supplies the \emph{kinetic} type -- a filled
window is photon-/neutrino-like, not a massive triple. There the gap was named: the \emph{derived}
rule computing the kinetic part, via the duality, as a $\xi$-fractal fraction of the static one is
missing. This document takes up exactly that gap. It is the \enquote{further work} on the dynamic
sector announced in Doc.~285.

% ==============================================================================
\section{The fixed foundation: the reversible Hilbert-space basis}
% ==============================================================================
Before developing the dynamic sector, let the fixed foundation it sits on be recorded. The FFGFT QM
translation (Doc.~282) is closed in three pieces:
\begin{enumerate}
  \item \textbf{The unambiguous Hilbert map} (Type~III, Doc.~270): $\mathcal H=L^2(T^4)\otimes
  \mathbb C^3$, states $=$ fiber-valued fields, QM operation $=$ unitary parallel transport,
  entanglement $=$ holonomy around cycles (Doc.~230), three masses $\leftrightarrow(M,r,\theta)$. A
  \emph{bijection} -- fully reversible.
  \item \textbf{The time-free mass form:} mass sits on a separate tensor factor; CHSH $=2\sqrt2$
  stays identical with and without the $\mathbb C^3$ fiber to $4\cdot10^{-16}$ -- no static QM result
  changes. Time is \emph{not computed} here; the reduced dynamics is time-local, unitary, Markovian.
  \item \textbf{The time form:} adds exactly \emph{one} effect -- non-Markovian memory in the reduced
  dynamics (backflow $0\to0.53$, CP-divisibility violated) -- which itself still lies within QM: a
  traced-out unitary bath reproduces it (Doc.~282/283). At the full level, therefore, unitary/reversible.
\end{enumerate}
Reversibility is precise \emph{per type} (Doc.~270): Type~III (Hilbert) fully reversible; Type~I
(mass circle $\to$ time, \enquote{time taken into the space}) \emph{almost} reversible -- lossless in
mode content, only the winding number stays ambiguous; Type~II (spatial projection $D3\!\to\!D0$)
\emph{not} reversible. The only genuine irreversibility is the spatial projection -- the measurement /
decompactification (Doc.~270/284) --, not the time-inclusion. The non-Markovian memory is not a break
of reversibility but the reduced view.

\paragraph{Aim of the dynamic sector.} On this fixed, reversible foundation, the actual task is the
\emph{Hilbert-space extension and its translation}: \emph{another} description \emph{out of the Hilbert space} -- the complementary,
decompactified/projected side that maps computationally, via $D_f=3-\xi$ and $\xi$, onto SI units. The
SI map is exactly the measurement/projection side (Type~I/II), and $D_f$ and $\xi$ are the bridge
quantities (Doc.~270/249/133). The dynamic sector is thus not a departure from the foundation but its
unrolling into measurable units -- and the sought kinetic $\xi$-fraction is one building block of
exactly that unrolling.

% ==============================================================================
\section{What is fixed}
% ==============================================================================
The duality fixes the \emph{placement}, not the size. With $\tilde T=1/\max(m,\omega)$
(Doc.~067, 080), an almost-massless window state has $\omega\gg m$, hence $\tilde T=1/\omega$: the
state sits in the \emph{kinetic} sector. That is secure. So is the conceptual separation:
\begin{center}
\begin{tabular}{lll}
\toprule
 & static & kinetic \\
\midrule
form        & $E_0=mc^2$ (rest)            & $E=hf$ (free) \\
seat        & Z$_3$ core (winding)         & window, when a dynamics runs \\
time        & $\tilde T=1/m$               & $\tilde T=1/\omega$ \\
status      & derived (Doc.~285)           & \emph{open} (this document) \\
\bottomrule
\end{tabular}
\end{center}
A distinction between two kinetic quantities matters: the ordinary kinetic energy of the electron,
$(\gamma-1)mc^2$ (Doc.~080), is \emph{not} $\xi$-suppressed. The $\xi$-fractal fraction sought here
is the \emph{window rest-mass} of an otherwise massless state -- not the kinetic energy of a massive
one.

% ==============================================================================
\section{The single data point: the neutrino}
% ==============================================================================
Doc.~047 treats the neutrino as a \enquote{near-massless photon}: $E^2=(pc)^2+(\sqrt{\xi^2/2}\,
mc^2)^2$, with $v_\nu=c\,(1-\xi^2/2)\approx0.9999999911\,c$ and $m_\nu\sim\tfrac{\xi^2}{2}m_e$.
This is the only existing value for a window rest-mass as a $\xi$-power of a core mass. It is,
however, explicitly an \emph{ansatz} (from the velocity/photon analogy), not \emph{derived} from the
duality. Doc.~204 additionally supplies a $\xi$-correction to the kinetic energy ($+0.69\,\xi_0$).
There are no further data points; in particular, a single value is not enough for the $\Delta m^2$
spectrum.

% ==============================================================================
\section{Narrowing the $\xi$-power: $p=2$}
% ==============================================================================
With the ansatz $m_{\text{window}}=C\cdot\xi^p\cdot m_{\text{ref}}$, tested against the bounds
(reproducing the $\xi^2/2$ value; the photon limit; dimensional consistency; an oscillation floor
$\gtrsim\!\sqrt{\Delta m^2_{\text{atm}}}\approx 49$~meV; a direct kinematic bound $m_\nu\lesssim0.8$~eV (KATRIN, model-independent)),
only $p=2$ is compatible with the meV scale: $p=1$ gives $\mathcal{O}(10\text{--}70)$~eV (far above
the bound), $p\ge3$ falls below every anchor. So $p=2$ is the \emph{unique} order-of-magnitude
compatible power -- a genuine order-of-magnitude bound, not a fit.

\paragraph{Numerology warning.} \emph{Within} $p=2$, the coefficient ($\tfrac12$) and the reference
mass are \emph{not} fixed by the single anchor. Different references land somewhere in the neutrino
window (e.g.\ $\xi^2/2\cdot m_e\approx4.5$~meV, $\xi^2\cdot m_e\approx9$~meV,
$\xi^2/2\cdot\sqrt{m_e m_\mu}\approx65$~meV). This is \emph{not} evidence but the P35 trap:
coefficient and reference need a derivation, not tuning. Computed: \texttt{kinetik\_xi\_anteil.py}.

% ==============================================================================
\section{The structural lead: $m_{\text{window}}=\xi\cdot\omega$}
% ==============================================================================
The Doc.~047 velocity expression $v_\nu=c(1-\xi^2/2)$ is exactly equivalent to a $\xi^1$ relation.
For $m\ll p$ one has $v/c=1-\tfrac12(m/p)^2$; substituting gives $(m/p)^2=\xi^2$, hence
\begin{equation}
m_{\text{window}}=\xi\cdot p\;\;\big(\approx \xi\cdot\omega\ \text{via}\ \tilde T\cdot m=1\big).
\end{equation}
The \enquote{$\xi^2$} in the notation \enquote{mass $=\xi^2/2\cdot m_e$} is thus presumably the
\emph{square image} of a fundamental $\xi^1$ relation: the static rest-mass is $\xi$ times the
kinetic scale. This is the natural anchor for a derivation \emph{via the duality}, because
$\omega=1/\tilde T$ is exactly the duality quantity. The open question therefore shifts from
\enquote{why $\xi^2$?} to \enquote{why $\xi^1$ between rest-mass and kinetic scale, and what sets
$\omega$ absolutely?}.

% ==============================================================================
\section{What remains to be derived}
% ==============================================================================
Three points are explicitly open:
\begin{enumerate}
  \item \textbf{The reference/scale -- largely closed} (next section). The two lab anchors select the
  reference as the \emph{single} FFGFT scale $E_0=\sqrt{m_e m_\mu}$; in its units $\omega/E_0=\xi/2$
  and $m_{\text{window}}/E_0=\xi^2/2$. Residual: the \emph{one} structural statement $\omega=\xi\,E_0$
  (exactly one fractal $\xi$-power).
  \item \textbf{The coefficient.} The factor $\tfrac12$ in $\xi^2/2$ is not derived; possibly it is
  purely the velocity square ($\tfrac12(m/p)^2$), i.e.\ no free factor at all.
  \item \textbf{The spectrum.} A single $\xi$-fractal value sets the \emph{scale}, not the
  \emph{spectrum} ($\Delta m^2_{\text{sol}}$, $\Delta m^2_{\text{atm}}$, mixing angles). That demands
  the specified $T^7$ dynamics (the second open computation, Doc.~285).
\end{enumerate}

% ==============================================================================
\section{Closing the scale question: the lab anchors select $E_0$}
% ==============================================================================
The first open point -- what sets $\omega$ absolutely -- closes further than a free parameter would. In
FFGFT there is only \emph{one} absolute scale; everything else is dimensionless ratios. The question is
therefore not \enquote{which number} but \enquote{which of the existing scales}. With
$m_{\text{window}}=\tfrac{\xi^2}{2}\cdot\text{reference}$, the two \emph{lab} anchors (oscillation floor
$\gtrsim49$~meV; KATRIN $\lesssim0.8$~eV) bracket the reference:
\begin{center}
\begin{tabular}{lccc}
\toprule
reference & $m_{\text{window}}$ & floor $\ge49$~meV & KATRIN $\le0.8$~eV \\
\midrule
$m_e$ & $4.5$~meV & \emph{fails} (too light) & yes \\
$E_0=\sqrt{m_e m_\mu}$ & $65$~meV & yes & yes \\
$M\approx314$~MeV & $2.8$~eV & yes & \emph{fails} \\
$m_\mu$ & $940$~meV & yes & \emph{fails} (too heavy) \\
\bottomrule
\end{tabular}
\end{center}
$m_e$ is too light, $M$ and $m_\mu$ too heavy -- \emph{only} $E_0=\sqrt{m_e m_\mu}$ fits, and that is
exactly the one natural FFGFT scale (Doc.~282, $\lambda_e\lambda_\mu=E_0$). The reference is thus not a
free quantity but the existing scale, selected by the lab.

In units of $E_0$ everything then becomes a pure $\xi$-power times $\tfrac12$:
\begin{equation}
\frac{\omega}{E_0}=\frac{\xi}{2},\qquad \frac{m_{\text{window}}}{E_0}=\frac{\xi^2}{2}.
\end{equation}
No free parameter remains. The $\tfrac12$ is the velocity square ($\tfrac12(m/p)^2$, Doc.~047) --
kinematic, not a tuning factor.

\paragraph{What honestly remains.} The lab window ($49$--$800$~meV) is a factor $\sim$few wide: it
selects the \emph{scale} $E_0$ sharply (excluding the alternatives by orders of magnitude) but does not
pin the exact coefficient by itself. What remains is exactly \emph{one} structural statement,
\begin{equation}
\omega=\xi\,E_0\,,
\end{equation}
-- $E_0$ suppressed by exactly \emph{one} fractal $\xi$-power ($D_f=3-\xi$). This is no longer
\enquote{free scale, self-referential} but a non-circular computation step: $\omega$ as the
$\xi$-suppressed eigenvalue of the connection-Laplacian in the time-extended Hilbert structure
(Doc.~283 §8). It closes the point -- or falsifies it. Computed: \texttt{schliesse\_omega.py}.

% ==============================================================================
\section{The $\xi^1$ mechanism: lifting the zero mode by the dimension deficit}
% ==============================================================================
The remaining residual -- why $\omega=\xi\,E_0$ with exactly \emph{one} $\xi$-power -- has a concrete
mechanism. The connection-Laplacian on $T^4$ has an exact zero mode (the holonomy-flat constant). The
fractal correction $g=\eta(1+\xi f)$ acts as a diagonal potential perturbation and lifts it;
first-order perturbation theory gives
\begin{equation}
\delta\lambda_0 = \xi\,\langle 0|f|0\rangle = \xi\cdot\overline f\,,
\end{equation}
i.e.\ the power $\xi^1$ \emph{precisely when} the correction has a non-vanishing torus average
$\overline f\neq0$; for $\overline f=0$ the first order vanishes and the mode lifts only at $\xi^2$.

$D_f=3-\xi$ is a dimension \emph{deficit} -- a systematic, non-volume-preserving change of the measure
--, hence $\overline f\neq0$. So the power is \emph{one}. Numerically (\texttt{xi1\_eigenwert.py}): the
zero-mode lift scales with $p=0.999$ for a deficit ($\overline f\neq0$) versus $p=1.98$ for a
volume-preserving ripple ($\overline f=0$) -- so the power is not accidental but follows from the
nature of $D_f=3-\xi$.

\paragraph{Status.} The $\xi^1$ power is thus \emph{derived} (zero-mode lift by the dimension deficit),
not asserted. What remains is only the $O(1)$ coefficient -- kinematically the $\tfrac12$ (Doc.~047) --
and confirming that the physical kinetic state is this lifted zero mode. With that, $\omega=\tfrac{\xi}{2}E_0$
is closed up to an $O(1)$ factor.

% ==============================================================================
\section{Relation to the static sector}
% ==============================================================================
The static sector is closed and computed: the Z$_3$ mass operator ($r=\sqrt2$, $\theta=2/9$) plus the
mass circle yields Koide $2/3$, the lepton ratios and -- as a pre-registered, matched-null-tested
declaration -- the BD12-XCORE interface (Doc.~285). The dynamic sector is the \emph{complementary,
open} side: the same parameter $\xi$, but as a \emph{fraction} rather than an \emph{eigenvalue}. A
derivation of $m_{\text{window}}=\xi\cdot\omega$ would move the window-filling dynamics (Doc.~263,
holographic scale projection) from description to computation, and the neutrino-like $3'$ candidate
(Doc.~285) from conditional to derived.

% ==============================================================================
\section{Reproducibility}\label{sec:repro-286}
% ==============================================================================
\texttt{kinetik\_xi\_anteil.py} (numpy-only): narrowing of the $\xi$-power against the bounds, the
duality framing ($\tilde T=1/\max(m,\omega)$), and the finding $p=2$ with the explicitly open
reference/coefficient/spectrum gap. The script claims no derivation; it narrows. \texttt{schliesse\_omega.py}: the lab selection of the reference $E_0$ (table) and the reduction $\omega/E_0=\xi/2$, $m_{\text{window}}/E_0=\xi^2/2$, with the $\xi^1$ remainder as the marked closing point.
